Dans l'espace rapporté à un repère orthonormé \Oijk, on considère les points :

\begin{center}
A$(4~;~-4~;~4)$, \quad B$(5~;~-3~;~2)$, \quad C$(6~;~- 2;~3)$, \quad D(5~;~1~;~1)
\end{center}

\begin{enumerate}
\item Démontrer que le triangle ABC est rectangle en $B$.
\item Justifier qu’une équation cartésienne du plan (ABC) est :

\[
x - y - 8 = 0.
\]

\item On note $d$ la droite passant par le point $D$ et orthogonale au plan (ABC).

\medskip

	\begin{enumerate}
		\item Déterminer une représentation paramétrique de la droite $d$.
		\item On note H le projeté orthogonal du point $D$ sur le plan $(ABC)$.
		
Déterminer les coordonnées du point H.
		\item Montrer que DH $= 2\sqrt 2$.
	\end{enumerate}
\item 
	\begin{enumerate}
		\item Montrer que le volume de la pyramide ABCD est égal à $2$.

On rappelle que le volume V d'une pyramide se calcule à l'aide de la formule :
\[V = \dfrac13 \times \mathcal{B} \times h\]
où $\mathcal{B}$ est l'aire d'une base de la pyramide et $h$ la hauteur correspondante.
		\item On admet que l'aire du triangle BCD est égale à $\dfrac{\sqrt{42}}{2}$.

En déduire la valeur exacte de la distance du point A au plan (BCD).
	\end{enumerate}
\end{enumerate}
